\documentclass[12pt, a4paper]{article} 
\usepackage[utf8]{inputenc}
\usepackage[english]{babel}
\usepackage{graphicx}
\usepackage{amsmath,amssymb}
\usepackage{lmodern}
\usepackage{tikz}
\usepackage{tabularx}
\usepackage{makecell}
\usetikzlibrary{arrows.meta, positioning}
\usepackage{iftex}
\ifPDFTeX
  \usepackage[T1]{fontenc}
  \usepackage[utf8]{inputenc}
  \usepackage{textcomp} % provide euro and other symbols
\else % if luatex or xetex
  \usepackage{unicode-math}
  \defaultfontfeatures{Scale=MatchLowercase}
  \defaultfontfeatures[\rmfamily]{Ligatures=TeX,Scale=1}
\fi
\usepackage{amsmath, amssymb, amsfonts}
\usepackage[left=2.5cm, right=2.5cm, top=2.5cm, bottom=2.5cm]{geometry}
\usepackage{caption} 
\usepackage{booktabs}
\usepackage{hyperref}
\hypersetup{
    colorlinks=true,
    linkcolor=blue,
    filecolor=magenta,
    urlcolor=cyan,
}
\usepackage{fancyhdr}
\usepackage{enumitem} 
\usepackage{lipsum}
\usepackage{float}
\usetikzlibrary{positioning, arrows.meta}

\pagestyle{fancy}
\fancyhf{} 
\fancyhead[L]{\nouppercase{\rightmark}}
\fancyhead[R]{\thepage} 
\fancyfoot[C]{}
\renewcommand{\headrulewidth}{0.4pt} 
\renewcommand{\footrulewidth}{0pt} 

% Trang bia
\title{
    \textbf{Ho Chi Minh University of Science} \\
    \textbf{Department of Information Technology} \\
    \vspace{0.5cm}
    \includegraphics[width=0.5\textwidth]{img/logohcmus.png} \\
    \vspace{1cm}
    \huge \textbf{Lab02 Report: Hashiwokakero} \\
    \vspace{1cm}   
    \Large \textbf{Course: Introduction to Artificial Intelligence} \\
    \Large \textbf{Instructor: Vo Nhat Tan} \\
}
\author{
    \shortstack[l] {
        \textbf{23127296 - Nguyen Thanh Luan} \\
        \textbf{23127302 - Tran Quang Phuc} \\
        \textbf{23127539 - Nguyen Thanh Tien} \\
        \textbf{23127543 - Vu Van Vu} \\
    }
}
\date{July, 2025} 

\begin{document}

\maketitle 

% muc luc
\newpage 
\tableofcontents 

% introduction
\newpage 

\section{Introduction}
\begin{table}[htbp]
    \centering
    \caption{Contribution table}
    \label{tab:contribution}
    \begin{tabular}{@{}lc@{}} % bo khoang trong
        \toprule
        \textbf{Work} & \textbf{Member} \\
        \midrule
        Formulate CNF constraints & Nguyen Thanh Luan\\
        Automate CNF generation & Nguyen Thanh Luan \\
        Design input files & Tran Quang Phuc \\
        Implement backtracking & Nguyen Thanh Tien \\
        Implement brute-force & Vu Van Vu  \\
        Implement A* & Tran Quang Phuc \\
        Solve using PySat & Nguyen Thanh Luan \\
        \bottomrule
    \end{tabular}
\end{table}

\begin{table}[htbp]
    \centering
    \caption{Requirement completion table}
    \label{tab:contribution}
    \begin{tabular}{@{}lc@{}} % bo khoang trong
        \toprule
        \textbf{Requirements} & \textbf{Completion level} \\
        \midrule
        Design at least 10 input files & Yes\\
        Formulate CNF constraints & Yes \\
        Automate CNF generation & Yes \\
        Use PySat to solve & Yes\\
        Implement backtracking & Yes\\
        Implement brute-force & Yes\\
        Implement A* & Yes\\
        Detailed explanation of algorithms & Yes \\
        Results and evaluation & Yes \\
        Report & Yes\\
        \bottomrule
    \end{tabular}
\end{table}
\section{Problem Introduction}
Hashiwokakero (also known as Bridges) is a Japanese puzzle. It is played on a rectangular grid. There are islands on the grid. The player can draw bridges between the islands in a straight line. The goal is to make all the islands form a single connected group. The puzzle follows the following rules:
\begin{itemize}
    \item They must begin and end at distinct islands, travelling a straight line in between.
    \item They must not cross any other bridges or islands.
    \item They may only run orthogonally (i.e. they may not run diagonally).
    \item At most two bridges connect a pair of islands.
    \item The number of bridges connected to each island must match the number on that island.
    \item The bridges must connect the islands into a single connected group.
\end{itemize}
\section{Encode Hashiwokakero to CNF}
\subsection{Logical Variable Definition}
To obtain the CNF formula for the Hashiwokakero puzzle, first we introduce a logical variable \\
\[ x_{i,j,i',j',k} \]
where 
\begin{itemize}
  \item $(i,j)$ and $(i',j')$ is the location of the 2 islands.
  \item $k$ is the number of bridges between 2 islands.($k=1,2$)
\end{itemize}
\subsection{CNF constraints}
The following constraints are converted to CNF in order to solve the Hashiwokakero puzzle:
\subsubsection{Bridges must begin and end at distinct islands, travelling in a straight line}
\begin{itemize}
    \item The endpoints $(i,j)$ and $(i',j')$ must be different: $(i,j) \ne (i',j')$.
    \item The bridge must lie on the same row or the same column:
    \[
    \forall i, j, i', j', k \quad \left( x_{i,j,i',j',k} \rightarrow \left[ (i = i' \land j \ne j') \lor (j = j' \land i \ne i') \right] \right)
    \]
    
\end{itemize}
\subsubsection{Bridges must not cross any other bridges}
Two bridges $x_{i_1,j_1,i_1',j_1',k_1}$ and $x_{i_2,j_2,i_2',j_2',k_2}$ intersect if:

\begin{itemize}
  \item $i_1 = i_1'$, $j_2 = j_2'$ (one runs horizontal, one runs vertical),
  \item $i_2 < i_1 < i_2'$, and $j_1 < j_2 < j_1'$ (they cross each other),
  \item $\{i_1,j_1\}, \{i_1',j_1'\} \ne \{i_2,j_2\}, \{i_2',j_2'\}$ (they don't run from same islands).
\end{itemize}
\subsubsection{At most two bridges connect a pair of islands}
The number of bridges is equal or smaller than 2:
\[
    \forall i, j, i', j', k \quad ( x_{i,j,i',j',k} \rightarrow 0 \leq k \leq 2)
    \]
\subsubsection{The number of bridges connected to each island must match the number on that island}
Let $B_{i,j}$ be the number of bridges connecting to the island located at ($i$, $j$).

\[
\sum_{\substack{i',j',k \\ (i',j',k) \in D}} k \cdot x_{i,j,i',j',k} = B_{i,j}
\]

where 
\begin{itemize}
    \item $D$ is the set of all islands $(i',j')$ that can go directly to $(i,j)$.
\end{itemize}
\subsection{Implementation}
\subsubsection{Generating bridge variables}
Firstly, the logical variables are generated by looping through the input file. The idea is to consider 4 direction from an island ($i$, $j$). If an island ($i'$, $j'$) is met during a direction, then two variables $x_{i,j,i',j',1}$ and $x_{i,j,i',j',2}$ are added to the list of variables.
\subsubsection{Adding constraints ensuring that $x_{i,j,i',j',1}$ and $x_{i,j,i',j',2}$ cannot coexist}
For every pair $x_{i,j,i',j',1}$ and $x_{i,j,i',j',2}$, we add the following condition to ensure they cannot coexist.
\[
\lnot v_1 \lor \lnot v_2
\]
\subsubsection{Adding constraints ensuring that the number of bridges to the island must match the number on that island}
For every island, we do the following:
\begin{itemize}
    \item Get all the logical variables representing bridges connect directly to that island.
    \item Find the combinations of those logical variables that satisfy the condition.
    \item We define a new variable $c$ for every combination.
    \item For every combination {$v_1$, $v_2$, ..., $v_n$}, we add \(\lnot v_1 \lor \lnot v_2 \lor \dots \lor v_n \lor c\), and for every logical variable $v_i$ in the combination, we add \(\lnot c \lor v_i\). These steps are the equivalence of adding \(c \iff v_1 \land v_2 \land \dots \land v_n\).
    \item For every combination, we also add the negative of all the logical variables in first step that are not in the combinations.
    \item For every pair of combinations $c_1$ and $c_2$, we add \(\lnot c_1 \lor \lnot c_2\), to ensure at most one combination is valid.
    \item For all combinations {$c_1$, $c_2$, ..., $c_n$}, we add \(\ c_1 \lor \ c_2 \lor \dots \lor c_n\) to ensure at least one combination is correct.
\end{itemize}
\subsubsection{Adding constraints ensuring no bridges cross each other}
For every logical variable $x_{i_1,j_1,i_1',j_1',k_1}$ (or $v_1$), we do the following:
\begin{itemize}
    \item For every logical $x_{i_2,j_2,i_2',j_2',k_2}$ (or $v_2$), check the following condition:
    \begin{itemize}
        \item $i_1 = i_1'$, $j_2 = j_2'$ (one runs horizontal, one runs vertical),
        \item $i_2 < i_1 < i_2'$, and $j_1 < j_2 < j_1'$ (they cross each other),
        \item $\{i_1,j_1\}, \{i_1',j_1'\} \ne \{i_2,j_2\}, \{i_2',j_2'\}$ (they don't run from same islands).
    \end{itemize}
    \item If all the conditions are satisfied we add \(\lnot v_1 \lor \lnot v_2\) to ensure these bridges cannot coexist.
\end{itemize}

\section{Algorithm Explanation}

\subsection{Backtracking}

\subsubsection{Intuition}
Backtracking is a trial-and-error technique. If at some point the current path cannot lead to a valid solution, the algorithm backtracks.

\subsubsection{Implementation}
\begin{enumerate}
  \item \textbf{Generate variables:} Assign Boolean variables to bridges between islands.

  \item \textbf{Add logical constraints:}
  \begin{itemize}
    \item Ensure no double connections 
    \item Enforce the degree constraints.
    \item Prevent bridges from crossing each other.
  \end{itemize}

  \item \textbf{Initialize the assignment:} 
  Create an assignment list to track which variables are currently set to \texttt{True} (bridge exists), \texttt{False} (bridge does not exist), or \texttt{None} (not yet assigned).

  \item \textbf{Partial constraint checking:} 
  During backtracking, after assigning a value to a variable, check whether all fully-assigned clauses are still satisfied. This prevents exploring invalid branches.

  \item \textbf{Recursive backtracking:}
  Try assigning \texttt{True} or \texttt{False} to each variable in order.
  \begin{itemize}
    \item If the partial assignment does not violate any constraints, proceed to assign the next variable.
    \item If a violation is detected, undo the assignment and try the opposite value.
    \item If all variables are assigned, check if the resulting graph is fully connected.
  \end{itemize}

  \item \textbf{Solution or failure:}
  If a complete valid assignment is found and the graph is connected, the solution is displayed.
  Otherwise, the algorithm concludes that no valid solution exists.
\end{enumerate}
\subsection{Brute-force}

\subsubsection{Intuition}
The brute-force approach tries all possible combination of bridges to find the solution.
It guarantees correctness, but consumes great amount of time.

\subsubsection{Implementation}

\begin{enumerate}
  \item \textbf{Generate variables:} Define Boolean variables for all bridges.

  \item \textbf{Add logical constraints:}
  \begin{itemize}
    \item Prevent double bridges.
    \item Ensure that the total number of bridges connected to each island matches its degree.
    \item Forbid crossing bridges.
  \end{itemize}

  \item \textbf{Iterate through all assignments:}
  Use exhaustive enumeration to try every possible combination of \texttt{True} or \texttt{False} for all variables.

  \item \textbf{Constraint checking:}
  For each assignment, check whether all logical clauses are satisfied. If any clause is unsatisfied, skip to the next assignment.

  \item \textbf{Check connectivity:}
  If the assignment satisfies all clauses, check whether the resulting graph is connected.

  \item \textbf{Display result or fail:}
  If a valid combination is found, display the solution and stop. If after all assignments but there's no solution -> this puzzle is unsolvable.
\end{enumerate}
\subsection{A*}

\subsubsection{Intuition}
A* is a best-first search algorithm that explores the most promising partial assignments first, guided by a cost function.

\subsubsection{Heuristics}
To guide the search efficiently, A* uses a heuristic function $h(n)$ that estimates the "cost to goal" from a given partial assignment. The total cost function is defined as:
\[
f(n) = g(n) + h(n)
\]
where:
\begin{itemize}
  \item $g(n)$ is the actual number of variables assigned so far (the path cost).
  \item $h(n)$ is the estimated number of unsatisfied clauses.
\end{itemize}

Several heuristics can be used:

\begin{itemize}
  \item \textbf{Clause-based heuristic:} Count the number of clauses that are currently unsatisfied under the partial assignment.

  \item \textbf{Optimistic heuristic:} Count only clauses that are fully assigned and violated. This makes the heuristic admissible and allows for more exploration.

  \item \textbf{Graph-based penalty (optional):} Add a penalty for partial solutions that look disconnected or under-connected, so the algorithm can be more likely to have globally connected configurations.
\end{itemize}

\subsubsection{Implementation}

\begin{enumerate}
  \item \textbf{Generate variables:} Represent all possible bridges between islands using Boolean variables.

  \item \textbf{Add constraints:} Encode logical conditions into CNF clauses:
  \begin{itemize}
    \item No duplicate bridges.
    \item Degree of each island must match its number.
    \item No crossing bridges.
  \end{itemize}

  \item \textbf{Initialize priority queue:} Start with an empty assignment and compute its $f(n)$ value.

  \item \textbf{Search loop:}
  \begin{enumerate}
    \item Pop the assignment with the smallest $f(n)$ from the priority queue.
    \item If all variables are assigned and all clauses are satisfied, and the resulting graph is connected, return the solution.
    \item Otherwise, pick one unassigned variable and expand the current state by assigning it both \texttt{True} and \texttt{False}.
    \item For each expanded state, compute $f(n)$ and push it into the priority queue.
    \item Optionally skip states that already violate a clause under full assignment.
  \end{enumerate}

  \item \textbf{Fail case:} If the priority queue is exhausted without finding a solution, report failure.
\end{enumerate}


\section{Results}
\subsection{Test cases}
\subsection{Results}
\subsection{Algorithm Comparisons}

\end{document}